%==============================================================
%  lecons/algebre/complexes.tex — Nombres complexes
%==============================================================

\lecontitle{Nombres complexes}{Algèbre — Structure, géométrie et analyse}

%=============================================================
%  § 1 — STRUCTURE ALGÉBRIQUE
%=============================================================
\secnum{navyblue}{1}{Structure algébrique de $\mathbb{C}$}

\begin{tcolorbox}[thmbox]
  \textbf{Définition.}\enspace\index{nombres complexes}
  Le corps des nombres complexes est
  \[
    \mathbb{C} = \mathbb{R}[X]/(X^2+1),
  \]
  c'est-à-dire l'anneau des polynômes réels modulo $X^2+1$.
  On note $i$ la classe de $X$, vérifiant
  \[
    \boxed{i^2 = -1.}
  \]
  Tout élément s'écrit de façon unique $z = a + ib$ avec $a,b\in\mathbb{R}$.
  On pose $\mathrm{Re}(z)=a$ (\emph{partie réelle}\index{partie réelle}) et
  $\mathrm{Im}(z)=b$ (\emph{partie imaginaire}\index{partie imaginaire}).

  \medskip
  \textbf{Opérations.}
  \begin{itemize}
    \item Addition : $(a+ib)+(c+id) = (a+c)+i(b+d)$.
    \item Multiplication : $(a+ib)(c+id) = (ac-bd)+i(ad+bc)$.
    \item Inverse : $(a+ib)^{-1} = \dfrac{a-ib}{a^2+b^2}$ pour $(a,b)\neq(0,0)$.
  \end{itemize}

  \medskip
  \textbf{Conjugué et module.}\enspace\index{conjugué}\index{module}
  Pour $z = a+ib$ on définit :
  \[
    \bar z = a - ib \quad\text{(conjugué)},\qquad
    |z| = \sqrt{a^2+b^2} \quad\text{(module)}.
  \]
\end{tcolorbox}

\rubriq{Propriétés algébriques fondamentales}

\begin{mdframed}[style=proofbar]
  Pour tous $z, w \in \mathbb{C}$ :
  \begin{enumerate}
    \item $\overline{z+w} = \bar z + \bar w$,\quad
          $\overline{zw} = \bar z\,\bar w$,\quad
          $\bar{\bar z} = z$.
    \item $z\bar z = |z|^2$, donc $z^{-1} = \bar z / |z|^2$ si $z\neq 0$.
    \item $|zw| = |z|\,|w|$,\quad $|z+w| \leq |z|+|w|$
          (\emph{inégalité triangulaire}\index{inégalité triangulaire}).
    \item $\mathrm{Re}(z) = \dfrac{z+\bar z}{2}$,\quad
          $\mathrm{Im}(z) = \dfrac{z-\bar z}{2i}$.
    \item $|\mathrm{Re}(z)| \leq |z|$,\quad $|\mathrm{Im}(z)| \leq |z|$.
  \end{enumerate}
\end{mdframed}

\rubriq{Représentation géométrique}
$\mathbb{C}$ s'identifie au plan $\mathbb{R}^2$ : $z = a+ib \leftrightarrow (a,b)$.
Le module $|z|$ est la distance à l'origine, et $|z-w|$ la distance entre $z$ et $w$.
Le conjugué $\bar z$ est le symétrique de $z$ par rapport à l'axe réel.

\decsep{}

%=============================================================
%  § 2 — FORME POLAIRE ET EXPONENTIELLE
%=============================================================
\secnum{navyblue}{2}{Forme polaire et formule d'Euler}

\begin{tcolorbox}[thmbox]
  \textbf{Forme polaire.}\enspace\index{forme polaire}
  Tout $z \in \mathbb{C}^*$ s'écrit
  \[
    z = r\,e^{i\theta}, \qquad r = |z| > 0,\quad \theta \in \mathbb{R}.
  \]
  $\theta$ est un \emph{argument}\index{argument} de $z$, défini modulo $2\pi$.
  L'argument principal $\arg(z) \in (-\pi,\pi]$ est l'unique représentant dans cet intervalle.

  \medskip
  \textbf{Formule d'Euler.}\enspace\index{formule d'Euler}
  \[
    \boxed{e^{i\theta} = \cos\theta + i\sin\theta.}
  \]
\end{tcolorbox}

\rubriq{Preuve de la formule d'Euler}

\begin{mdframed}[style=proofbar]
  On définit $\exp(z) = \sum_{k=0}^\infty z^k/k!$ sur $\mathbb{C}$ (série absolument convergente).
  En séparant termes pairs et impairs :
  \[
    e^{i\theta}
    = \sum_{k=0}^\infty \frac{(i\theta)^k}{k!}
    = \sum_{m=0}^\infty \frac{(-1)^m\theta^{2m}}{(2m)!}
      + i\sum_{m=0}^\infty \frac{(-1)^m\theta^{2m+1}}{(2m+1)!}
    = \cos\theta + i\sin\theta.
    \hfill\blacksquare
  \]
\end{mdframed}

\rubriq{Propriétés de la forme polaire}
Pour $z = re^{i\theta}$ et $w = se^{i\varphi}$ :
\[
  zw = rs\,e^{i(\theta+\varphi)}, \qquad
  \frac{z}{w} = \frac{r}{s}\,e^{i(\theta-\varphi)}, \qquad
  \bar z = re^{-i\theta}, \qquad
  \arg(zw) = \arg(z)+\arg(w) \pmod{2\pi}.
\]

\rubriq{Cercle unité}
L'ensemble $\mathbb{U} = \{z\in\mathbb{C} : |z|=1\} = \{e^{i\theta}:\theta\in\mathbb{R}\}$
est un groupe multiplicatif. Les formules de trigonométrie s'en déduisent immédiatement :
\[
  \cos\theta = \mathrm{Re}(e^{i\theta}) = \frac{e^{i\theta}+e^{-i\theta}}{2}, \qquad
  \sin\theta = \mathrm{Im}(e^{i\theta}) = \frac{e^{i\theta}-e^{-i\theta}}{2i}.
\]

\decsep{}

%=============================================================
%  § 3 — RÉSULTATS FONDAMENTAUX
%=============================================================
\secnum{navyblue}{3}{Résultats fondamentaux}

\begin{tcolorbox}[thmbox]
  \textbf{Théorème de De Moivre.}\enspace\index{De Moivre (théorème de)}
  Pour tout $\theta\in\mathbb{R}$ et tout $n\in\mathbb{Z}$ :
  \[
    (\cos\theta + i\sin\theta)^n = \cos(n\theta) + i\sin(n\theta).
  \]
\end{tcolorbox}

\begin{mdframed}[style=proofbar]
  Immédiat : $(e^{i\theta})^n = e^{in\theta}$. \hfill$\blacksquare$
\end{mdframed}

\rubriq{Linéarisation}
Le binôme de Newton appliqué à $(e^{i\theta}+e^{-i\theta})^n = (2\cos\theta)^n$
permet d'exprimer $\cos^n\theta$ en somme de $\cos(k\theta)$.
De même $\sin^n\theta$ via $(e^{i\theta}-e^{-i\theta})^n = (2i\sin\theta)^n$.

\begin{mdframed}[style=exbar]
  \textbf{Exemple — linéarisation de $\cos^3\theta$.}
  \[
    2^3\cos^3\theta = (e^{i\theta}+e^{-i\theta})^3
    = e^{3i\theta}+3e^{i\theta}+3e^{-i\theta}+e^{-3i\theta}
    = 2\cos(3\theta)+6\cos\theta.
  \]
  Donc $\cos^3\theta = \tfrac{1}{4}\cos(3\theta)+\tfrac{3}{4}\cos\theta$.
\end{mdframed}

\begin{tcolorbox}[thmbox]
  \textbf{Racines $n$-ièmes de l'unité.}\enspace\index{racines de l'unité}
  Pour $n\geq 1$, l'équation $z^n = 1$ admet exactement $n$ solutions dans $\mathbb{C}$ :
  \[
    \mathbb{U}_n = \left\{\omega_k = e^{2i\pi k/n} : k = 0, 1, \ldots, n-1\right\}.
  \]
  Ces racines forment un groupe cyclique pour la multiplication, engendré par
  $\omega = e^{2i\pi/n}$.
  Géométriquement, ce sont les sommets d'un polygone régulier à $n$ côtés
  inscrit dans le cercle unité.

  \medskip
  \textbf{Somme et produit.}
  \[
    \sum_{k=0}^{n-1} \omega^k = 0 \quad (n\geq 2), \qquad
    \prod_{k=0}^{n-1} \omega^k = (-1)^{n+1}.
  \]
\end{tcolorbox}

\begin{mdframed}[style=proofbar]
  $z^n - 1 = \prod_{k=0}^{n-1}(z-\omega^k)$ par factorisation.
  La somme des racines est le coefficient de $z^{n-1}$ (nul) et le produit
  est le terme constant $(-1)^n\cdot(-1)^n\cdot 1 = (-1)^{n+1}$.\hfill$\blacksquare$
\end{mdframed}

\begin{tcolorbox}[thmbox]
  \textbf{Racines $n$-ièmes d'un complexe.}\enspace\index{racines $n$-ièmes}
  Pour $a = \rho\,e^{i\alpha} \in \mathbb{C}^*$, les solutions de $z^n = a$ sont :
  \[
    z_k = \rho^{1/n}\,e^{i(\alpha + 2k\pi)/n}, \qquad k = 0, 1, \ldots, n-1.
  \]
  Il y a exactement $n$ racines, régulièrement espacées sur un cercle de rayon $\rho^{1/n}$.
\end{tcolorbox}

\begin{tcolorbox}[thmbox]
  \textbf{Théorème fondamental de l'algèbre (d'Alembert--Gauss).}\enspace
  \index{théorème fondamental de l'algèbre}
  Tout polynôme non constant à coefficients complexes admet au moins une racine dans $\mathbb{C}$.
  En conséquence, tout polynôme de degré $n\geq 1$ se factorise en $n$ facteurs linéaires sur $\mathbb{C}$ :
  \[
    P(z) = a_n\prod_{k=1}^{n}(z - z_k), \qquad z_1,\ldots,z_n \in \mathbb{C}.
  \]
  Autrement dit : $\mathbb{C}$ est \emph{algébriquement clos}\index{algébriquement clos}.
\end{tcolorbox}

\rubriq{Conséquence pour les polynômes réels}
Si $P\in\mathbb{R}[X]$ et $z_0$ est une racine, alors $\bar z_0$ l'est aussi.
Les racines complexes non réelles arrivent donc par paires conjuguées,
et $P$ se factorise sur $\mathbb{R}$ en facteurs de degré 1 et 2.

\decsep{}

%=============================================================
%  § 4 — EXEMPLES, CONTRE-EXEMPLES, EXERCICES
%=============================================================
\secnum{exgreen}{4}{Exemples, pièges et exercices}

\noindent{\bfseries\color{navyblue} Exemples.}

\begin{mdframed}[style=exbar]
  \textbf{1.\ Résolution de $z^4 = -16$.}\enspace
  $-16 = 16\,e^{i\pi}$, donc les 4 racines sont
  $z_k = 2\,e^{i(\pi+2k\pi)/4} = 2\,e^{i\pi(2k+1)/4}$, $k=0,1,2,3$ :
  \[
    z_0 = \sqrt{2}(1+i),\quad z_1 = \sqrt{2}(-1+i),\quad
    z_2 = \sqrt{2}(-1-i),\quad z_3 = \sqrt{2}(1-i).
  \]

  \smallskip\noindent
  \textbf{2.\ Formule de duplication.}\enspace
  De $e^{2i\theta} = (e^{i\theta})^2 = (\cos\theta+i\sin\theta)^2$ on tire :
  \[
    \cos(2\theta) = \cos^2\theta - \sin^2\theta, \qquad
    \sin(2\theta) = 2\sin\theta\cos\theta.
  \]

  \smallskip\noindent
  \textbf{3.\ Racines cubiques de l'unité.}\enspace
  $j = e^{2i\pi/3} = -\tfrac{1}{2}+i\tfrac{\sqrt{3}}{2}$ vérifie
  $j^3=1$, $1+j+j^2=0$, $j^2 = \bar j = e^{-2i\pi/3}$.
\end{mdframed}

\vspace{6pt}
\noindent{\bfseries\color{counterred} Pièges.}

\begin{mdframed}[style=cexbar]
  \textbf{$\sqrt{z\cdot w} \neq \sqrt{z}\cdot\sqrt{w}$ sur $\mathbb{C}$.}\enspace
  Par exemple $\sqrt{(-1)(-1)} = \sqrt{1} = 1$ mais
  $\sqrt{-1}\cdot\sqrt{-1} = i\cdot i = -1$.
  Les identités algébriques sur $\sqrt{\phantom{x}}$ supposent des réels positifs.

  \smallskip\noindent
  \textbf{$\arg$ n'est pas un homomorphisme.}\enspace
  La règle $\arg(zw)=\arg(z)+\arg(w)$ n'est vraie que modulo $2\pi$.
  Si l'on travaille avec l'argument principal $\in(-\pi,\pi]$, un saut de $\pm 2\pi$
  peut apparaître et rompre des égalités naïves.

  \smallskip\noindent
  \textbf{$\overline{f(z)} \neq f(\bar z)$ en général.}\enspace
  C'est vrai pour $f\in\mathbb{R}[X]$ (polynôme à coefficients réels), mais pas
  pour une fonction quelconque. Par exemple $f(z)=iz$ donne
  $\overline{f(z)} = -i\bar z \neq i\bar z = f(\bar z)$.
\end{mdframed}

\vspace{6pt}
\exolabel

\begin{mdframed}[style=exobar]
  \begin{enumerate}
    \item Mettre $z = \dfrac{(1+i)^{10}}{(1-i\sqrt{3})^5}$ sous forme polaire puis algébrique.

    \item Résoudre dans $\mathbb{C}$ : $z^2 - (3+i)z + (2+3i) = 0$.

    \item Montrer que pour tout $n\geq 2$ :
          $\displaystyle\sum_{k=0}^{n-1} e^{2ik\pi/n} = 0$
          en calculant la somme géométrique.

    \item Linéariser $\sin^5\theta$ en somme de $\sin(k\theta)$.

    \item Soit $P\in\mathbb{R}[X]$ de degré impair. Montrer sans utiliser le
          théorème fondamental de l'algèbre que $P$ admet au moins une racine réelle.

    \item Calculer $\displaystyle\sum_{k=0}^{n-1}\cos(k\theta)$ et
          $\displaystyle\sum_{k=0}^{n-1}\sin(k\theta)$ pour $\theta \notin 2\pi\mathbb{Z}$.
  \end{enumerate}
\end{mdframed}

\leconfooter{R.\ Descartes (1637) — L.\ Euler (1748, $e^{i\pi}+1=0$) — C.\ F.\ Gauss (1799, th.\ fondamental)}
